\section{Conclusion}
\label{sec:conclusion}

This paper establishes a systematic methodology for understanding the performance and GPU-energy behavior of non-disaggregated and disaggregated LLM serving. We compare Native, prefill--decode (PD), and attention--FFN (AF) architectures in a common SGLang implementation under matched models, request streams, SLOs, and accelerator budgets. The study spans dense and MoE models, diverse input/output regimes, DP/TP/EP and hybrid parallelism, datacenter and consumer GPUs, NVLink/PCIe/RDMA communication, compute--memory capability, and GPU operating settings. By measuring latency, throughput, power, and total energy over the same request interval, the methodology avoids treating lower power as synonymous with greater energy efficiency.

The central premise of this work is that serving architecture cannot be evaluated independently of workload, placement, communication, and hardware. Disaggregation offers specialization and independent scaling, but its benefits must offset state transfer, activation communication, synchronization, and additional device residency. Likewise, a throughput-optimal parallel plan or a lower-frequency operating point is not necessarily energy optimal. Reproducible deployment manifests, communication validation, SLO-aware filtering, per-GPU energy accounting, and repeated open-loop trials provide the evidence needed to identify these trade-offs.

Our characterization is intended to support both system designers and operators: it provides a common basis for selecting architectures and deployment configurations and motivates future controllers that jointly optimize architecture, parallelism, placement, and frequency. As LLM inference becomes a growing component of datacenter electricity demand, such architecture-aware performance--energy analysis will be essential for building scalable and sustainable serving systems.
